\section{Introduction}

% OUTLINE ONLY — each bullet becomes a paragraph.

\begin{itemize}
  \item Hook: the compiler is inside every deployed model's trust boundary, and
    nobody audits it for confidentiality --- it can \emph{introduce} leaks the
    source never had, and (surprise finding) \emph{remove} leaks the source did
    have.
  \item The gap: side-channel work targets hand-written crypto kernels; verified
    compilation targets correctness, not secrecy; DL compilers
    (Inductor, MLIR pipelines, Polygeist) have neither.
  \item Two-sided thesis of the paper:
    \begin{itemize}
      \item \textbf{Measure} --- a differential eager-vs-compiled methodology
        attributes each observed channel to source or compiler (the
        $2{\times}2$,
        \texttt{docs/research/leak\_check.methodology.siddharth.md}).
      \item \textbf{Verify} --- per-pass checked specifications where a lowering
        is safe only if it both computes the same thing \emph{and} preserves
        non-interference (\texttt{prototypes/fcvd\_ct}).
    \end{itemize}
  \item Contributions list (draft):
    \begin{enumerate}
      \item A unified threat model: observer rings O0--O7 $\times$ leaked-asset
        axis, composing four independently developed frameworks
        (\texttt{docs/research/threat-models.agents.md}).
      \item The differential measurement rig and activation-function corpus,
        with the \texttt{exp} channel \emph{compiler-removed} as headline
        (\texttt{prototypes/leak\_check}).
      \item Two-property translation validation (``safe = congruent +
        non-interfering'') instantiated on Polygeist's eight-step pipeline,
        with SMT semantics for the memref$\rightarrow$LLVM slice
        (\texttt{prototypes/fcvd\_ct}).
      \item Negative-space honesty: the count-channel confound, retired probes,
        and model-blocked steps reported as first-class results.
    \end{enumerate}
  \item Roadmap paragraph.
\end{itemize}
